% GENERATED FILE. Edit canonical lessons, then run npm run generate:books.
% canonical-source-hash: fnv1a64:713b68b8a9013abd
% canonical-lessons: UR-C518-avaz, UR-C518-surila, UR-C518-rag, UR-C518-riyaz, UR-C518-gungunana

\chapter{Learning to Sing}
\label{ch:ur-learning-to-sing}
\hlchaptermodality{518}

\begin{chapteropening}
\textbf{By the end of this chapter:} \emph{I can talk about voices, ragas and practice.}

Complete the last lesson of chapter 518: I can talk about voices, ragas and practice.
\end{chapteropening}

\section[\texorpdfstring{\ur{آواز} (āvāz)}{آواز (āvāz)}]{\ur{آواز} (āvāz) — a voice; a sound}

\label{lesson:UR-C518-avaz}



\begin{quote}
Before the new one: say the Urdu for a dialogue, then the Urdu for to charm.
\end{quote}

\subsection*{You'll want to know: \ur{آواز}}
\textbf{\ur{آواز}} — \emph{āvāz} — \textquotedblleft{}a voice; a sound\textquotedblright{}.

\begin{quote}
\ur{اس} \ur{گلوکار} \ur{کی} \ur{آواز} \ur{میٹھی} \ur{ہے۔}
\end{quote}

\begin{quote}
\emph{is gulūkār kī āvāz mīṭhī hai.}
\end{quote}

\begin{quote}
This singer has a sweet voice.
\end{quote}

Feminine.

\begin{cousinweb}
\textbf{\ur{آواز}} (\emph{āvāz}) is usually traced to Persian \emph{āvāz}, \textquotedblleft{}a voice, a sound, a song\textquotedblright{}.
\end{cousinweb}

\begin{grammarlens}[title={What you've built}]
One more word for describing what you lived and felt.
\end{grammarlens}

\subsection*{Your turn}
\begin{itemize}
\raggedright
  \item \emph{Say it:} \emph{āvāz}
  \item \emph{Say it:} \emph{āvāz}, once more
  \item \emph{Say it:} \emph{is gulūkār kī āvāz mīṭhī hai.} — This singer has a sweet voice.
  \item \emph{Recall:} say the Urdu for a dialogue, then the Urdu for to charm, then say \emph{āvāz} again
\end{itemize}

\begin{tcolorbox}[breakable,colback=teal!4,colframe=teal!35!black,title={Before you move on}]
What does \ur{آواز} mean? (\textquotedblleft{}A voice; a sound\textquotedblright{}.) Say it once more.
\end{tcolorbox}
\section[\texorpdfstring{\ur{سریلا} (surīlā)}{سریلا (surīlā)}]{\ur{سریلا} (surīlā) — melodious; in tune}

\label{lesson:UR-C518-surila}



\begin{quote}
Before the new one: say the Urdu for to charm, then the Urdu for a voice.
\end{quote}

\subsection*{You'll want to know: \ur{سریلا}}
\textbf{\ur{سریلا}} — \emph{surīlā} — \textquotedblleft{}melodious; in tune\textquotedblright{}.

\begin{quote}
\ur{سارا} \ur{کی} \ur{آواز} \ur{سریلی} \ur{ہے۔}
\end{quote}

\begin{quote}
\emph{Sārā kī āvāz surīlī hai.}
\end{quote}

\begin{quote}
Sara has a lovely singing voice.
\end{quote}

\begin{cousinweb}
\textbf{\ur{سریلا}} (\emph{surīlā}) is inherited: CDIAL lists the Hindi cousin \emph{surīlā}, \textquotedblleft{}melodious\textquotedblright{}, under Sanskrit \emph{sura}, \textquotedblleft{}sound\textquotedblright{}, a variant of \emph{svara} (CDIAL 13498).
\end{cousinweb}

\begin{grammarlens}[title={What you've built}]
One more word for describing what you lived and felt.
\end{grammarlens}

\subsection*{Your turn}
\begin{itemize}
\raggedright
  \item \emph{Say it:} \emph{surīlā}
  \item \emph{Say it:} \emph{surīlā}, once more
  \item \emph{Say it:} \emph{Sārā kī āvāz surīlī hai.} — Sara has a lovely singing voice.
  \item \emph{Recall:} say the Urdu for to charm, then the Urdu for a voice, then say \emph{surīlā} again
\end{itemize}

\begin{tcolorbox}[breakable,colback=teal!4,colframe=teal!35!black,title={Before you move on}]
What does \ur{سریلا} mean? (\textquotedblleft{}Melodious; in tune\textquotedblright{}.) Say it once more.
\end{tcolorbox}
\section[\texorpdfstring{\ur{راگ} (rāg)}{راگ (rāg)}]{\ur{راگ} (rāg) — a raga (a melodic mode)}

\label{lesson:UR-C518-rag}



\begin{quote}
Before the new one: say the Urdu for a voice, then the Urdu for melodious.
\end{quote}

\subsection*{You'll want to know: \ur{راگ}}
\textbf{\ur{راگ}} — \emph{rāg} — \textquotedblleft{}a raga (a melodic mode)\textquotedblright{}.

\begin{quote}
\ur{استاد} \ur{نے} \ur{مجھے} \ur{شام} \ur{کا} \ur{ایک} \ur{راگ} \ur{سکھایا۔}
\end{quote}

\begin{quote}
\emph{ustād ne mujhe shām kā ek rāg sikhāyā.}
\end{quote}

\begin{quote}
My teacher taught me an evening raga.
\end{quote}

Masculine.

\begin{cousinweb}
\textbf{\ur{راگ}} (\emph{rāg}) is inherited from Sanskrit \emph{rāga}, \textquotedblleft{}colouring; passion; a musical mode\textquotedblright{} (Monier-Williams).
\end{cousinweb}

\begin{grammarlens}[title={What you've built}]
One more word for describing what you lived and felt.
\end{grammarlens}

\subsection*{Your turn}
\begin{itemize}
\raggedright
  \item \emph{Say it:} \emph{rāg}
  \item \emph{Say it:} \emph{rāg}, once more
  \item \emph{Say it:} \emph{ustād ne mujhe shām kā ek rāg sikhāyā.} — My teacher taught me an evening raga.
  \item \emph{Recall:} say the Urdu for a voice, then the Urdu for melodious, then say \emph{rāg} again
\end{itemize}

\begin{tcolorbox}[breakable,colback=teal!4,colframe=teal!35!black,title={Before you move on}]
What does \ur{راگ} mean? (\textquotedblleft{}A raga (a melodic mode)\textquotedblright{}.) Say it once more.
\end{tcolorbox}
\section[\texorpdfstring{\ur{ریاض} (riyāz)}{ریاض (riyāz)}]{\ur{ریاض} (riyāz) — practice (in music)}

\label{lesson:UR-C518-riyaz}



\begin{quote}
Before the new one: say the Urdu for melodious, then the Urdu for a raga.
\end{quote}

\subsection*{You'll want to know: \ur{ریاض}}
\textbf{\ur{ریاض}} — \emph{riyāz} — \textquotedblleft{}practice (in music)\textquotedblright{}.

\begin{quote}
\ur{سارا} \ur{ہر} \ur{صبح} \ur{دو} \ur{گھنٹے} \ur{ریاض} \ur{کرتی} \ur{ہے۔}
\end{quote}

\begin{quote}
\emph{Sārā har subah do ghanṭe riyāz kartī hai.}
\end{quote}

\begin{quote}
Sara practises for two hours every morning.
\end{quote}

Masculine.

\begin{cousinweb}
\textbf{\ur{ریاض}} (\emph{riyāz}) is usually traced, through Persian, to Arabic \emph{riyāḍa}, \textquotedblleft{}exercise, training\textquotedblright{}; Urdu uses it for a musician's daily practice.
\end{cousinweb}

\begin{grammarlens}[title={What you've built}]
One more word for describing what you lived and felt.
\end{grammarlens}

\subsection*{Your turn}
\begin{itemize}
\raggedright
  \item \emph{Say it:} \emph{riyāz}
  \item \emph{Say it:} \emph{riyāz}, once more
  \item \emph{Say it:} \emph{Sārā har subah do ghanṭe riyāz kartī hai.} — Sara practises for two hours every morning.
  \item \emph{Recall:} say the Urdu for melodious, then the Urdu for a raga, then say \emph{riyāz} again
\end{itemize}

\begin{tcolorbox}[breakable,colback=teal!4,colframe=teal!35!black,title={Before you move on}]
What does \ur{ریاض} mean? (\textquotedblleft{}Practice (in music)\textquotedblright{}.) Say it once more.
\end{tcolorbox}
\section[\texorpdfstring{\ur{گنگنانا} (gungunānā)}{گنگنانا (gungunānā)}]{\ur{گنگنانا} (gungunānā) — to hum}

\label{lesson:UR-C518-gungunana}



\begin{quote}
Before the new one: say the Urdu for a raga, then the Urdu for practice (in music).
\end{quote}

\subsection*{You'll want to know: \ur{گنگنانا}}
\textbf{\ur{گنگنانا}} — \emph{gungunānā} — \textquotedblleft{}to hum\textquotedblright{}.

\begin{quote}
\ur{سارا} \ur{باورچی} \ur{خانے} \ur{میں} \ur{ایک} \ur{گیت} \ur{گنگنا} \ur{رہی} \ur{تھی۔}
\end{quote}

\begin{quote}
\emph{Sārā bāvarchī khāne meṅ ek gīt gungunā rahī thī.}
\end{quote}

\begin{quote}
Sara was humming a song in the kitchen.
\end{quote}

Past, with \ur{نے} (he or she …): \ur{اس} \ur{نے} \ur{ایک} \ur{گیت} \ur{گنگنایا}, \emph{us ne ek gīt gungunāyā}.

\begin{cousinweb}
\textbf{\ur{گنگنانا}} (\emph{gungunānā}) is inherited: CDIAL lists the Hindi cousin \emph{gungunānā}, \textquotedblleft{}to murmur\textquotedblright{}, under \emph{*ganagana}, a form imitating the sound (CDIAL 4013).
\end{cousinweb}

\begin{grammarlens}[title={What you've built}]
One more word for describing what you lived and felt.
\end{grammarlens}

\subsection*{Your turn}
\begin{itemize}
\raggedright
  \item \emph{Say it:} \emph{gungunānā}
  \item \emph{Say it:} \emph{gungunānā}, once more
  \item \emph{Say it:} \emph{Sārā bāvarchī khāne meṅ ek gīt gungunā rahī thī.} — Sara was humming a song in the kitchen.
  \item \emph{Recall:} say the Urdu for a raga, then the Urdu for practice (in music), then say \emph{gungunānā} again
\end{itemize}

\begin{tcolorbox}[breakable,colback=teal!4,colframe=teal!35!black,title={Before you move on}]
What does \ur{گنگنانا} mean? (\textquotedblleft{}To hum\textquotedblright{}.) Say it once more.
\end{tcolorbox}
